\medskip



On considère le point A(1~;~1~;~0) et le vecteur $\vect{u}\begin{pmatrix}0\\2\\- 1\end{pmatrix}$.

On considère le plan $\mathcal{P}$ d'équation : $x + 4y + 2z + 1 = 0$.

\medskip

\begin{enumerate}
\item 


On sait que $M(x~;~y~;~z) \in (d) \iff \vect{\text{A}M} = t \vect{u}$, avec $t \in \R$.

Soit $M(v) \in (d) \iff \left\{\begin{array}{l c l}
x - 1&=&0t\\
y - 1&=&2t\\
z - 0&=&-t
\end{array}\right., \: t \in \R \iff \left\{\begin{array}{l c l}
x &=&1\\
y&=&1 + 2t\\
z&=&\phantom{1 + } -t
\end{array}\right., \: t \in \R$.
\item 
Un point de coordonnées $(x~;~y~;~z)$ est commun à la droite $(d)$ et au plan $\mathcal{P}$ si ces coordonnées les équations de la droite et du plan soit le système :

$\left\{\begin{array}{l c l}
x &=&1\\
y&=&1 + 2t\\
z&=&\phantom{1 + } -t\\
x + 4y + 2z + 1&=&0
\end{array}\right., \: t \in \R$.

En remplaçant $x, \:y$ et $z$ par leurs valeurs en fonction de $t$ dans la dernière équation on obtient :

$1 + 4(1 + 2t) + 2\times (- t) + 1 = 0 \iff 1 + 4 + 8t - 2t + 1 = 0 \iff 6t + 6 = 0 \iff t = - 1$

En remplaçant $t$ par $- 1$ dans les trois équations de $(d)$, on obtient :

$\left\{\begin{array}{l c l}
x &=&1\\
y&=&1 - 2\\
z&=&\phantom{1 + } 1
\end{array}\right.$

On a donc B$(1~;~- 1~;~1) \in (d) \cap \mathcal{P}$.
\item 
	\begin{enumerate}
		\item 
		
		On a $\vect{\text{AB}}\begin{pmatrix}0\\-2\\1\end{pmatrix}$ et $\vect{\text{AC}}\begin{pmatrix}0\\-2\\- 1\end{pmatrix}$ : ces deux vecteurs ne sont pas colinéaires puisque les deux premières coordonnées sont égales mais pas les troisièmes. ; les points A, B et C non alignés définissent bien le plan (ABC).
		\item 
		On a $\vect{n} \cdot \vect{\text{AB}} = 0 + 0 + 0 = 0$ et 
		
$\vect{n} \cdot \vect{\text{AC}} = 0 + 0 + 0 = 0$.

Conclusion : le vecteur $\vect{n}$ est orthogonal à deux vecteurs non colinéaires du plan (ABC) : il est donc normal à ce plan.
		\item 
On sait qu'alors $M(x~;~y~;~z) \in (\text{ABC}) \iff x \times 1 + y\times 0 + z \times 0 + d = 0$, avec $d \in \R$.

Or A$(1~;~1~;~0) \in (\text{ABC}) \iff 1 \times 1 + d = 0  \iff d = - 1$.

On a donc $M(x~;~y~;~z) \in (\text{ABC}) \iff x - 1 = 0$.
	\end{enumerate}
\item 
	\begin{enumerate}
		\item 
AB$^2 = \left\|\vect{\text{AB}}\right\|^2 = 0^2 + (- 2)^2 + 1^2 = 4 + 1 = 5$ ;

AC$^2 = \left\|\vect{\text{AB}}\right\|^2 = 0^2 + (- 2)^2 + (- 1)^2 = 4 + 1 = 5$.

On a donc AB$^2 = \text{AC}^2 \implies $ AB = AC : le triangle ABC est isocèle en A.
		\item 
		On a H$\left(\dfrac{1+1}{2}~;~\dfrac{(- 1) + (- 1)}{2}~;~\dfrac{1 - 1}{2} \right)$ soit H$(1~;~- 1~;~0)$.
		
On a donc AH$^2 = (1 - 1)^2 + (- 1 - 1)^2 + (0 - 0)^2 = 4$, d'où AH $ = 2$.

ABC étant isocèle en A, la droite (AH) médiane est aussi hauteur relative au côté [BC].

O a donc $\mathcal{A}(\text{ABC}) = \dfrac{\text{BC} \times \text{AH}}{2}$.

Comme BC$^2 = (1 - 1)^2 + (- 1 + 1)^2 + (- 1 - 1)^2 = 0 + 0 + 4$, d'où BC $ = 2$.

Donc $\mathcal{A}(\text{ABC}) = \dfrac{2 \times 2}{2} = 2$.
		

	\end{enumerate}
\item Soit D le point de coordonnées $(0~;~-1~;~1)$.
	\begin{enumerate}
		\item 
		Avec $\vect{\text{BD}}\begin{pmatrix}- 1\\0\\0\end{pmatrix}$, donc $\vect{\text{BD}} = - \vect{n}$ qui est un vecteur normal au plan (ABC), donc $\vect{\text{BD}}$ aussi et par conséquent la droite (BD) est une hauteur de la pyramide ABCD.
		\item 
		Pour la pyramide ABCD on prend comme base le triangle (ABC) et comme hauteur le segment [BD], d'où :
		
		$\mathcal{V}(\text{ABCD}) = \dfrac{\mathcal{A}\text{ABC}) \times [\text{BD}]}{3}$.
		
		Avec BD$^2 = \left\|\vect{\text{BD}}\right\|^2 = (- 1)^2 + 0^2 + 0^2 = 1$, on a donc BD $ = 1$ et 
		
		$\mathcal{V}(\text{ABCD}) = \dfrac{2\times 1}{3} = \dfrac23$ en unité d'aire.
	\end{enumerate}
	




\end{enumerate}

\bigskip
